% Date: June 13, 2026
% Status: R2 — Second Major Revision Cover Letter
%==============================================================================

\documentclass[11pt,a4paper]{letter}
\usepackage[utf8]{inputenc}
\usepackage[T1]{fontenc}
\usepackage{amsmath,siunitx}
\usepackage[margin=2.5cm]{geometry}
\usepackage{xr}
\externaldocument{Manuscript_JPCL_26-06-13}
\externaldocument{SI_JPCL_26-06-13}

\address{
Department of Physics\\
Faculty of Science\\
University of Douala, Cameroon\\
Email: steve.teguia@facsciences-uy1.cm
}

\signature{Steve Cabrel Teguia Kouam}

\begin{document}

\begin{letter}{
\textbf{Prof. Gregory Scholes}\\
Editor-in-Chief\\
The Journal of Physical Chemistry Letters\\
American Chemical Society
}

\opening{Dear Prof. Scholes,}

We submit the second major revision of our manuscript ``Selective Vibronic Excitation for Coherent Energy Transport in Photosynthetic and Agrivoltaic Systems'' (Manuscript ID: jz-2026-00994t.R1) for consideration as a Letter in \textit{The Journal of Physical Chemistry Letters}.

We thank both reviewers for their continued, careful evaluation of our work. The revisions address all points raised in the second-round reports (10 June 2026):

\begin{enumerate}
\item \textbf{Transfer yield redefined (Reviewer 3).} The forward transfer yield $\Phi_{\mathrm{FT}}$ is now defined as the long-time equilibrium population of the reaction-center-coupled site (BChl~3), replacing the previous survival-probability definition. The corrected enhancement $\eta = 0.18 \pm 0.04$ ($p < 0.001$) is physically transparent and directly reflects directed energy flow toward the reaction center.

\item \textbf{Vibronic-basis terminology clarified (Reviewer 2).} The term ``polaron-dressed states'' has been replaced throughout the main text by ``vibronically dressed states,'' consistent with the established literature (Womick \& Moran, JPC Lett.\ 2015; \v{S}anda et al., JPC Lett.\ 2022). The displaced-oscillator (polaron) frame argument is now confined to the SI derivation.

\item \textbf{Early vibronic-resonance literature cited (Reviewer 2).} The introduction now explicitly cites and discusses the 2015 and 2022 JPC Letters references on vibronic resonance in light-harvesting systems.

\item \textbf{Spectral density corrected (Reviewer 3).} Figure 3e now displays $J(\omega)$ over the full $0$--$2000$~cm$^{-1}$ range, correctly showing both the Drude--Lorentz background and all 12 discrete vibronic peaks. The reorganization energy ($\lambda_{\mathrm{tot}} = \SI{53}{\per\centi\meter}$) is justified with reference to the Kleinekath\"ofer/Coker parameterization and a new sensitivity check at $\lambda_D = \SI{100}{\per\centi\meter}$ (SI Section~S6).

\item \textbf{Dissipation confirmed (Reviewer 3).} The production ensemble was re-run with the verified canonical parameters ($\lambda_D = \SI{35}{\per\centi\meter}$). The revised Figure 3 clearly shows population decay toward thermal equilibrium and a well-defined decaying envelope in $C_{l_1}(t)$.

\item \textbf{Computational jargon removed (Reviewer 3).} All hardware-specific language has been removed from the main text. Figure~2 (the resource-architecture flowchart) has been moved to SI Section~S5. The manuscript now contains only physics-motivated statements on the methodology.

\item \textbf{PT-HOPS scaling and Matsubara convergence clarified (Reviewer 3).} The $\mathcal{O}(1)$ size-invariance claim has been corrected; the Drude--Lorentz vs.\ discrete-mode Matsubara distinction is now explained explicitly.

\item \textbf{Proofreading completed (Reviewer 2).} The manuscript has been independently reviewed by two co-authors and an automated spell-checker; all identified typographical errors have been corrected.
\end{enumerate}

These updates reinforce our finding that selective vibronic excitation extends excitonic coherence lifetimes by \SIrange{20}{50}{\percent} and enhances directed transfer to the reaction center by $\eta = 0.18$ at room temperature. We believe the revisions fully address all reviewer concerns and that the manuscript now meets the standards of \textit{The Journal of Physical Chemistry Letters}.

\textbf{Response to manuscript formatting request (28-Apr-2026).}
We have addressed all non-scientific formatting items as follows:
\begin{enumerate}
\item \textbf{Section headings removed.} The prohibited section headings (Introduction, Results and Discussion, Conclusion) have been removed. Major divisions are now indicated by bold run-in paragraph leads (e.g., \textbf{Theoretical Framework.}), and subsections use unnumbered \texttt{\textbackslash subsection*\{\}} headings, consistent with the JPCL Letter format.
\item \textbf{TOC Graphic included.} A new TOC Graphic is included via the \texttt{tocentry} environment in the \texttt{achemso} class, immediately preceding the document body, as required.
\item \textbf{Supporting Information.} The Supporting Information file has been checked and contains no highlighting, tracked changes, or colored text.
\end{enumerate}

\textbf{Cover art (Supplementary cover invitation, 28-Apr-2026).}
We are pleased to submit the TOC Graphic (\texttt{TOC\_Graphic.png}) for consideration as a supplementary journal cover. The image depicts the dual-band spectral filtering scheme applied to the FMO complex, with the engineered photon bath selectively populating vibronic resonances. We consent to both front and supplementary cover consideration.

All authors have approved this revised manuscript. We declare no conflicts of interest.
We confirm this work received no external funding; validated ORCID iDs for the authors will be provided during the ACS submission. Top-level section headings were converted to bold run-in paragraph leads to comply with JPCL format.

Thank you for your continued consideration of our work.

\closing{Sincerely,}

\fromsig{
Steve Cabrel Teguia Kouam\\
Department of Physics\\
University of Douala, Cameroon\\
Email: steve.teguia@facsciences-uy1.cm
}

\end{letter}
\end{document}